%% =====================================================================
%% 多代理人立場交換辯論機制應用於投資決策
%% 編譯：pdfLaTeX + BibTeX（Overleaf：Menu → Compiler → pdfLaTeX）
%% 需同時上傳 references.bib
%% =====================================================================

\documentclass[12pt,a4paper]{article}

%% ---------- 套件 ----------
\usepackage[T1]{fontenc}
\usepackage{newtxtext}                  % 英文：Times
\usepackage{amsmath}
\usepackage[left=2.5cm,right=2.5cm,top=2.5cm,bottom=2.5cm]{geometry}
\usepackage{CJKutf8}                    % 中文：楷體（bkai）
\usepackage{graphicx,array,tabularx,booktabs,float}
\usepackage{setspace,enumitem,footmisc}
\usepackage[explicit]{titlesec}
\usepackage{caption}
\usepackage{natbib}
\usepackage{tikz}
\usetikzlibrary{positioning,arrows.meta,calc}
\usepackage[bookmarks=false,hidelinks]{hyperref}

%% ---------- 行距與段落 ----------
\setstretch{1.5}                        % 內文 1.5 倍行距
\setlength{\parindent}{2em}             % 首行縮排兩字
\setlength{\parskip}{0pt}
\setlist{itemsep=2pt,topsep=4pt,leftmargin=2.5em}
\renewcommand{\arraystretch}{1.3}
\setlength{\tabcolsep}{4.5pt}
\emergencystretch=2em
\setlength{\footnotesep}{8pt}

%% ---------- 中文數字編號 ----------
\newcommand{\chnum}[1]{\ifcase#1\or 一\or 二\or 三\or 四\or 五\or 六\or 七\or 八\or 九\or 十\fi}
\renewcommand{\thesection}{\chnum{\value{section}}}
\renewcommand{\thesubsection}{（\chnum{\value{subsection}}）}

%% ---------- 標題：大小與間距（全黑） ----------
\titleformat{\section}{\fontsize{16}{24}\selectfont}{\thesection、}{0em}{#1}
\titleformat{\subsection}{\fontsize{14}{21}\selectfont}{\thesubsection}{0em}{#1}
\titleformat{\paragraph}[runin]{\fontsize{12}{18}\selectfont}{}{0em}{#1}
\titleformat{name=\section,numberless}{\fontsize{16}{24}\selectfont}{}{0em}{#1}
\titlespacing*{\section}{0pt}{18pt}{12pt}
\titlespacing*{\subsection}{0pt}{12pt}{6pt}
\titlespacing*{\paragraph}{2em}{6pt}{0.5em}

%% ---------- 圖表說明 ----------
\captionsetup{font={small,stretch=1.2},labelsep=quad,skip=6pt}
\captionsetup[table]{position=top}

%% ---------- 頁碼 ----------
\pagestyle{plain}

%% ---------- 表格線條 ----------
\setlength{\heavyrulewidth}{0.8pt}
\setlength{\lightrulewidth}{0.4pt}

\newcommand{\grp}[1]{\textsf{#1}}

\newcolumntype{Y}{>{\raggedright\arraybackslash}X}

\begin{document}
\begin{CJK}{UTF8}{bkai}

\renewcommand{\refname}{參考文獻}
\renewcommand{\tablename}{表}
\renewcommand{\figurename}{圖}

%% =====================================================================
%% 封面
%% =====================================================================
\begin{titlepage}
\centering
\setstretch{1.6}
\vspace*{2cm}
{\fontsize{18}{27}\selectfont 國立陽明交通大學管理科學系\par}
\vspace{0.3cm}
{\fontsize{16}{24}\selectfont 專題研究報告\par}
\vspace{3.5cm}
{\fontsize{20}{32}\selectfont 多代理人立場交換辯論機制\\ 應用於投資決策\par}
\vspace{0.8cm}
{\fontsize{14}{21}\selectfont Multi-Agent Role-Switching Debate\\ for Investment Decision-Making\par}
\vfill
{\fontsize{14}{24}\selectfont
\begin{tabular}{r@{\hspace{1.5em}}l}
學\hspace{2em}生： & 黃岳振 \\
指導教授： & 林斯寅\ 教授 \\
\end{tabular}\par}
\vspace{2.5cm}
{\fontsize{14}{21}\selectfont 中華民國一一五年九月\par}
\vspace{1.5cm}
\end{titlepage}

%% =====================================================================
%% 摘要
%% =====================================================================
\pagenumbering{roman}
\begin{center}
{\fontsize{16}{24}\selectfont 摘\quad 要\par}
\end{center}
\vspace{0.3cm}

本專題以投資決策輔助為應用情境，設計並檢驗一套「立場交換辯論」的多代理人工作流。使用者輸入股票代號與分析日期後，Coordinator 設定資料時間邊界並分派任務，情緒面、基本面、技術面與總經代理人整理證據，合併為不含投資結論的中立研究報告；再由單次判斷、獨立投票、固定立場辯論與立場交換辯論四組決策機制讀取相同資料比較（獨立投票與兩辯論組之決策呼叫次數對齊），最後經守門員檢查並以 30 至 90 個交易日之回測評估。

多代理人辯論是否優於強單模型或模型集成仍有爭議 \citep{smit2023mad,zhang2025stop}；立場交換雖已用於問答與事實查核 \citep{liu2025roleswitch,chowdhury2026proclaim}，尚缺乏在無前視偏差、呼叫預算對齊之投資回測中的檢驗。因此本研究不預設立場交換有效，而在嚴格設計下檢驗之。本研究之立場交換以「自我反駁」為核心：交換輪中代理人改持相反立場、僅讀自己第一輪的論述，並須引用且反駁自己的原主張；RQ1 即比較此自我反駁機制（D）與一般固定立場辯論（C）。

正式 RQ1 以 qwen3:32b、季度協議 v3-1002.3 於 2026 年 10 月 2 日事前登記 180 個案例後執行：標的匿名化、買賣持有由模型直接決定、辯論組裁決者須擇一方向，並凍結程式原始碼與模型權重之雜湊。主要比較為 \grp{D} 對 \grp{C} 之 60 日方向準確率與 Sharpe 比率（兩者經 Holm 校正後皆顯著且為正方判定 \grp{D} 較佳）。180 案全數完成。\grp{D} 之方向準確率為 55.0\%，\grp{C} 為 58.3\%（差 $-3.3$ 個百分點，McNemar $p=0.21$）；年化 Sharpe 比率 \grp{D} 為 0.54、\grp{C} 為 0.59（差 $-0.05$，Ledoit--Wolf $p=0.80$），經 Holm 校正後皆不顯著，故不支持自我反駁機制優於固定立場辯論；此結論在 30／90 日、零成本、守門員後決策與排除低新聞密度案例下皆一致。四組均未勝過同期買進持有（Sharpe 1.01）。探索性分析中，以 Gemini 3 Flash 重跑同一批案例得到相同結論，且兩模型各組準確率皆低於一律看多（62.2\%），顯示限制可能來自輸入資訊。

\vspace{0.6cm}
\noindent 關鍵詞：大型語言模型、多代理人辯論、立場交換、投資決策、回測、前視偏差

\vspace{0.2cm}
\noindent Keywords: large language models, multi-agent debate, role switching, investment decision-making, backtesting, look-ahead bias

\vspace{0.2cm}
\noindent JEL 分類代碼：G11、G17、C45、C63\footnote{依 JEL 分類系統自行選定：G11（投資組合選擇與投資決策）、G17（金融預測與模擬）、C45（神經網路與相關方法）、C63（計算方法與模擬）。}

\clearpage
\begin{center}
{\fontsize{16}{24}\selectfont Abstract\par}
\end{center}
\vspace{0.3cm}

\noindent Using investment decision support as the application, this project designs and tests a multi-agent \emph{role-switching debate} workflow. After the user enters a ticker and an analysis date, the Coordinator sets the data time boundary and dispatches tasks; sentiment, fundamental, technical and macro agents organize the evidence, which is merged into a neutral research report that contains no investment conclusion. Four decision mechanisms---a single decision, independent voting, fixed-role debate and role-switching debate---then read the same data for comparison (independent voting and the two debate arms use the same number of decision calls). Finally, a Gatekeeper checks the outputs, and backtests over 30 to 90 trading days evaluate them.

Whether multi-agent debate outperforms a strong single model or a model ensemble remains contested \citep{smit2023mad,zhang2025stop}; role switching has been applied to question answering and fact-checking \citep{liu2025roleswitch,chowdhury2026proclaim}, but has not been tested in investment backtests that are free of look-ahead bias and matched in call budget. This study therefore does not presume that role switching is effective, and instead tests it under a strict design. The role switch here centers on self-rebuttal: in the switched round, each agent takes the opposite stance, reads only its own first-round argument, and must quote and rebut its own original claim. RQ1 compares this self-rebuttal mechanism (D) with ordinary fixed-role debate (C).

The formal RQ1 run used qwen3:32b under quarterly protocol v3-1002.3, with 180 cases preregistered on 2 October 2026: tickers were anonymized, Buy/Hold/Sell was decided directly by the model, debate adjudicators had to choose a direction, and hashes of the source code and model weights were frozen. The primary comparison is D versus C on 60-day direction accuracy and Sharpe ratio (D is judged better only if both are positive and significant after Holm adjustment). All 180 cases were completed. Direction accuracy was 55.0\% for D and 58.3\% for C (difference $-3.3$ percentage points, McNemar $p=0.21$); the annualized Sharpe ratio was 0.54 for D and 0.59 for C (difference $-0.05$, Ledoit--Wolf $p=0.80$). Neither is significant after Holm adjustment, so the results do not support the self-rebuttal mechanism outperforming fixed-role debate; this conclusion is consistent at 30 and 90 days, with zero costs, for post-Gatekeeper decisions and when low-news-density cases are excluded. No arm outperformed buy-and-hold over the same period (Sharpe 1.01). In an exploratory analysis, rerunning the same cases with Gemini 3 Flash led to the same conclusion, and every arm of both models had lower accuracy than always predicting a rise (62.2\%), suggesting that the limitation may lie in the input information.

\clearpage
\pagenumbering{arabic}

%% =====================================================================
\section{引言}
\label{sec:motivation}
%% =====================================================================

\subsection{研究背景}

大型語言模型已能閱讀新聞、財報、產業評論與市場資料，因此常被用來輔助投資分析，例如財報解讀、新聞情緒分析與交易決策 \citep{lopezlira2023chatgpt,araci2019finbert,liu2023fingpt}。但投資決策本身高度不確定：股價同時受公司基本面、技術走勢、市場情緒、總體經濟與突發事件影響。若只讓單一模型直接輸出買進、持有或賣出，模型可能因證據不足、敘事流暢而給出過度自信的建議。

因此，本研究把投資分析視為測試多代理人工作流的應用場景。重點不是證明 AI 能取代專業投資判斷，而是設計一個能把資料來源、正反論點、裁決理由與回測結果都留下紀錄的輔助流程，並以嚴謹的方式檢驗其中一個設計選擇是否真的有用。

\subsection{研究動機}

\paragraph{多代理人辯論的效果仍有爭議。}
\citet{du2023debate} 讓多個模型實例互相閱讀、批評並修正答案，顯示辯論可提升推理與事實正確性；\citet{liang2023divergent} 則指出模型自我反思容易固守既有答案，辯論可鼓勵發散思考。然而 \citet{smit2023mad} 比較多種辯論策略後發現，多代理人辯論不一定穩定優於自我一致性等其他方法，效果受設定影響；\citet{zhang2025stop} 進一步主張目前的評估常高估辯論，應重新思考評估方式並納入模型異質性。因此，「辯論比較準」不能被當作前提。

\paragraph{立場交換被提出以對抗固守立場，但尚未在投資情境檢驗。}
被指派立場的辯論者常持續強化原立場，而非真正回應對方證據。\citet{liu2025roleswitch} 因此在辯論中途讓雙方交換角色並回報不確定性；\citet{chowdhury2026proclaim} 在法庭式辯論框架中以角色交換作為一致性診斷；\citet{chen2026mpt} 以相反的人格觀點揭露並修正偏誤。這些研究皆屬問答、事實查核或偏誤緩解，並非投資決策，也未在呼叫預算對齊、無前視偏差的回測中檢驗。

\paragraph{金融場景的評估特別容易失真。}
LLM 的訓練資料可能涵蓋回測期間，使模型直接記得後續走勢 \citep{glasserman2023lookahead}；\citet{kong2026bias} 回顧 164 篇金融 LLM 論文後指出，前視偏差等偏誤常被忽略，使回測結果不可靠。辯論組通常也比單次判斷多用數倍的呼叫，若不對齊預算，優勢可能只是「多想幾次」。

\paragraph{本研究的做法。}
本研究延續立場交換的構想：第二輪中，看多者改為看空、看空者改為看多，且只能看到自己第一輪的論述，必須指出並反駁自己先前的主張；第三輪再換回原立場，整合換位時發現的反證。但不預設它有效，而在三個條件下檢驗：四組決策機制讀同一份鎖定的報告、獨立投票與兩辯論組之決策呼叫次數對齊、分析日之後的資料絕不進入輸入，且研究樣本於執行前事前登記。

\subsection{研究問題}

\begin{enumerate}[label=RQ\arabic*：,leftmargin=4em]
  \item 同一模型內之組間比較：在相同資料與相同決策呼叫次數下，自我反駁式立場交換辯論（\grp{D}；交換輪改持相反立場、僅讀自身第一輪並反駁之）之方向準確率與風險調整後報酬是否優於固定立場辯論（\grp{C}）？並以獨立投票（\grp{B}）與單次判斷（\grp{A}）比較其他決策流程。
  \item 跨模型比較：新版 RQ1 所得之四組結果與排序，在 qwen3:32b 與 Gemini 3 Flash 之間是否一致？跨模型案例需使用相同資料集 ID 與驗證政策，按股票與分析日配對，不合併兩模型的協議統計。
  \item 跨時間尺度比較：RQ1 所得之四組結果與排序，在以季為單位（每季一次分析日、60 個交易日預測期間）與以月為單位（每月一次分析日、20 個交易日預測期間）之設計下是否一致？
\end{enumerate}

\noindent RQ1 的核心比較為 \grp{D} 對 \grp{C}。\grp{D} 交換輪隔離對手論述是刻意設計：若同時看得到對手，代理人易轉為回應對手，而非檢討自己的主張。因此立場互換、隔離對手與強制引用自身原主張共同構成處理，\grp{D} 對 \grp{C} 估計的是此自我反駁機制相對一般辯論的效果。若要區分「換到相反立場」與「暫時不看對手」各自的貢獻，需另設消融組（例如僅讀自身歷史但不換立場），列為後續研究。\grp{B} 與 \grp{A} 比較其他決策流程。RQ2 與 RQ3 以新版季度結果為基準。

\subsection{研究貢獻}

\begin{enumerate}[label=\arabic*.]
  \item 將立場交換辯論應用於投資決策，並設計 \grp{B}、\grp{C}、\grp{D} 決策呼叫次數對齊、\grp{A} 為單次基準之四組流程比較；\grp{D} 以立場互換、隔離對手與強制引用自身原主張實現自我反駁，與一般辯論比較；不宣稱區分其中各要素的個別貢獻。本研究並非首次提出立場交換 \citep{liu2025roleswitch,chowdhury2026proclaim}，貢獻在於在金融情境中以明確界定的機制比較檢驗之。
  \item 建立具時間邊界、可追溯之多代理人回測框架：所有決策機制使用同一份不可變資料快照，模型輸出經逐項引用與數字驗證，研究樣本於執行前事前登記，每個案例可匯出完整稽核紀錄。
  \item 提供事前登記之實證：以 qwen3:32b 於 180 個季度案例檢驗自我反駁機制，未觀察到相對固定立場辯論之改善（方向準確率差 $-3.3$ 個百分點、Sharpe 差 $-0.05$，皆不顯著）。
\end{enumerate}

\subsection{研究範圍}

\begin{table}[H]
\centering\small
\caption{研究範圍}
\label{tab:scope}
\begin{tabularx}{\textwidth}{@{}l Y@{}}
\toprule
面向 & 範圍說明 \\
\midrule
專題定位 & 投資決策之 LLM 多代理人工作流與其公平比較，定位為決策輔助，不是自動交易系統 \\
決策輸出 & 買進、持有、賣出、不交易，以及預期報酬、信心、理由、引用證據與風險 \\
實作模組 & Coordinator、研究代理人、中立研究報告、決策機制層、短期工作記憶、長期決策記憶、裁決者、守門員、回測 \\
研究資料 & 九檔美國大型股（AAPL、NVDA、GOOGL、MSFT、AMZN、JPM、MCD、INTC、GE），2021 至 2025 年每季一個決策點，共 180 個案例 \\
\bottomrule
\end{tabularx}
\end{table}

%% =====================================================================
\section{文獻探討}
\label{sec:related}
%% =====================================================================

本章整理與本專題相關的四類文獻：多代理人辯論、角色交換與多視角推理、金融 LLM 與投資代理人，以及前視偏差與評估。目的不是完整比較所有方法，而是說明本專題的設計從哪些研究得到啟發，又為何不預設辯論有效。

\subsection{多代理人辯論}

\citet{du2023debate} 讓多個模型實例進行多輪辯論後收斂答案，改善推理與事實正確性；\citet{liang2023divergent} 指出模型自我反思容易固守既有答案，並以辯論鼓勵發散思考。另一方面，\citet{smit2023mad} 比較多種辯論策略，顯示多代理人辯論不一定穩定優於自我一致性等其他方法，且對設定敏感；\citet{zhang2025stop} 主張目前的評估常高估辯論，呼籲重新思考評估方式並採用異質模型。自我一致性 \citep{wang2023selfconsistency} 以多次取樣投票提升穩定性，是比較辯論效果時必要的對照組，本研究的獨立投票組即據此設計。由此，本研究把重點放在可追蹤的流程與公平比較（呼叫次數對齊），而不是預設辯論一定比較準。

\subsection{角色交換與多視角推理}

\citet{liu2025roleswitch} 提出不確定性感知的角色交換辯論：兩個辯論者先給出初始答案並互相詰問，再於辯論中途交換角色論證相反立場，然後各自回報信心與不確定性，最後由另一個裁判模型裁決；在 OpenBookQA 上，其消融實驗顯示角色交換與不確定性回報兩個階段均有貢獻。\citet{chowdhury2026proclaim} 的 PROCLAIM 以法庭角色（原告、辯護、法官）與漸進式檢索進行主張查核，並以角色交換作為一致性診斷；其消融結果顯示角色交換帶來約 4.2 個百分點的準確率提升，但整體 token 用量約為標準多代理人辯論的 11 倍。\citet{chen2026mpt} 的多人格思考以相反的人格與中立觀點反覆對話，使偏誤外顯並被修正。這些研究支持「從相反立場重新思考」的想法，但皆非投資決策，且第二項的效益是伴隨成本增加而來，凸顯呼叫預算對齊的必要。本研究因此採用看多、看空與裁決者的三角架構，並讓獨立投票與兩辯論組之決策呼叫次數對齊（單次判斷為基準）。

\subsection{金融 LLM 與投資代理人}

\citet{lopezlira2023chatgpt} 顯示 LLM 對金融新聞標題的判讀含有與後續報酬相關的資訊；FinBERT \citep{araci2019finbert} 與 FinGPT \citep{liu2023fingpt} 說明金融語言具有領域特性；FNSPID \citep{dong2024fnspid} 提供具時間戳記的金融新聞。TradingAgents \citep{xiao2024tradingagents} 與 StockAgent \citep{zhang2024stockagent} 把 LLM 代理人放入較接近真實投資流程的環境，強調分工、交易規則與回測。本研究的研究代理人分工與這些工作類似，並以呼叫預算對齊比較不同決策流程。

\subsection{前視偏差與評估}

\citet{glasserman2023lookahead} 指出以 GPT 產生之情緒訊號預測報酬時，模型可能透過對公司的既有知識引入前視偏差，並嘗試以匿名化緩解；\citet{kong2026bias} 歸納金融 LLM 常見的五種偏誤（前視、存活、敘事、目標與成本偏誤），並指出多數研究未明確處理。本研究因而在資料端嚴格排除分析日以後的資訊，在模型端以無資料回想測試檢視模型是否記得後續報酬，並在回測中扣除交易成本。

\subsection{文獻整理與本專題設計}

\begin{table}[H]
\centering\small
\caption{文獻與本專題設計對應}
\label{tab:litmap}
\begin{tabularx}{\textwidth}{@{}l Y Y@{}}
\toprule
文獻方向 & 對本專題的啟發 & 本專題採用方式 \\
\midrule
多代理人辯論 & 多個模型可互相檢查答案；但效果不穩定、易被高估 & 看多與看空辯論加裁決者；不預設有效，以對照組檢驗 \\
自我一致性 & 多次取樣投票可提升穩定性 & 獨立投票組，呼叫次數與辯論組對齊 \\
角色交換 & 從相反立場重新思考可降低固守立場 & 第二輪交換立場，且只能看到自己第一輪論述 \\
金融 LLM 與投資代理人 & 金融任務需要領域資料、分工與回測 & 四面向研究代理人、事前登記、扣除成本之回測 \\
前視偏差與評估 & 訓練資料可能涵蓋回測期間；評估須明確處理偏誤 & 資料時間邊界、無資料回想測試、證據與數字驗證 \\
\bottomrule
\end{tabularx}
\end{table}

%% =====================================================================
\section{研究資料}
\label{sec:data}
%% =====================================================================

研究範圍為九檔美國大型股（AAPL、\allowbreak NVDA、\allowbreak GOOGL、\allowbreak MSFT、\allowbreak AMZN、\allowbreak JPM、\allowbreak MCD、\allowbreak INTC、\allowbreak GE），分析日為 2021 至 2025 年共二十個季末，合計 180 個案例。每個案例於分析日建立一份資料快照，所有證據皆記錄可取得時間，僅分析日前已公開者得進入研究輸入。資料來源與時間處理方式如表~\ref{tab:data}。

新版資料共含 83,093 則新聞標題（每案中位數 235.5 則），180 案皆有新聞；其中 36 案僅有即時 API 取得之標題（約 50 則，集中於 2021--2023 年之 JPM、MCD、AMZN 等），MCD 2023-12-31 僅 2 則、未達產生情緒指標所需之 3 則。這些案例全數保留於主分析，並於事前登記之敏感性分析中排除後重算。

\begin{table}[H]
\centering\small
\caption{資料來源與時間邊界處理}
\label{tab:data}
\begin{tabularx}{\textwidth}{@{}l l Y@{}}
\toprule
面向 & 來源 & 內容與時間處理 \\
\midrule
技術面 & 調整後日行情 & 近 20 日報酬、20 日與 60 日均線相對位置、年化波動度；僅使用分析日前之價格。 \\
基本面 & SEC XBRL 財報 & 營收、淨利、營業現金流與前一年同期比較，以及負債對資產比率；依申報日期篩選，僅保留分析日前已公開者。 \\
情緒面 & 金融新聞標題（FNSPID、Alpha Vantage） & 分析日前 90 天內之全部新聞標題（不設則數上限），以固定版本之 FinBERT 逐則評分，目標公司與背景新聞分開彙整為平均分數與正負比例。 \\
總經面 & FRED／ALFRED & 使用歷史版本資料，僅採用分析日當時實際已發布之數值，排除事後修正。 \\
\bottomrule
\end{tabularx}
\end{table}

\subsection{變數定義}

反應變數為 60 日報酬（RET60）：於分析日後第一個交易日開盤進場，持有 60 個交易日後於開盤出場，並扣除 Corwin--Schultz 價差估計之交易成本 \citep{corwin2012spread}；30 與 90 個交易日作為穩健性檢查，方向準確率以 RET60 之正負判定。解釋資訊為表~\ref{tab:data} 之四個面向，整理為附有識別碼之證據，由研究代理人彙整為中立研究報告，供決策代理人引用。

\subsection{資料品質處理}

為確保快照僅反映分析日當時可得之資訊，本研究採取三項處理：其一，財報依申報日期而非會計期末篩選，並以相同會計科目、相同期間長度與前一年同期比較，避免混用不同期間；其二，年增率以前期絕對值為分母，使負數基期惡化時呈現負值，避免符號誤導；其三，快照以內容雜湊為識別碼，四組決策機制與所有重跑皆讀取同一份快照。

%% =====================================================================
\section{研究方法}
\label{sec:methods}
%% =====================================================================

\subsection{研究架構}

圖~\ref{fig:flow} 為自原始資料至報告結果之流程。協調者依股票與分析日派發任務；資料代理人建立快照；四個研究代理人分別整理一個面向，合併為中立研究報告後鎖定；四組決策機制讀取同一份報告；守門員檢查引用與風險門檻；回測模組計算實際報酬並寫入長期記憶，最後進行統計檢定。

\begin{figure}[H]
\centering
\begin{tikzpicture}[
  node distance=5mm and 7mm,
  box/.style={draw=black, align=center, font=\small, minimum height=9mm, inner sep=4pt, line width=0.5pt},
  arm/.style={box, fill=gray!12, minimum width=25mm},
  arr/.style={-{Stealth[length=2mm]}, line width=0.5pt}
]
\node[box] (in) {股票、分析日};
\node[box, right=of in] (snap) {資料快照\\（分析日前資料）};
\node[box, right=of snap] (res) {研究代理人\\中立研究報告};
\node[arm, below=11mm of snap, xshift=-44mm] (a) {\grp{A}\ 單次判斷};
\node[arm, right=3mm of a] (b) {\grp{B}\ 獨立投票};
\node[arm, right=3mm of b] (c) {\grp{C}\ 固定立場};
\node[arm, right=3mm of c] (d) {\grp{D}\ 立場交換};
\node[box, below=10mm of $(b)!0.5!(c)$] (gate) {守門員：引用驗證、資料覆蓋、風險門檻};
\node[box, below=of gate] (bt) {回測：開盤對開盤、扣除交易成本};
\node[box, below=of bt] (st) {統計檢定：配對檢定、Holm 校正、日期群聚敏感性};
\draw[arr] (in) -- (snap);
\draw[arr] (snap) -- (res);
\foreach \g in {a,b,c,d} \draw[arr] (res.south) -- (\g.north);
\foreach \g in {a,b,c,d} \draw[arr] (\g.south) -- (gate.north);
\draw[arr] (gate) -- (bt);
\draw[arr] (bt) -- (st);
\end{tikzpicture}
\caption{研究流程。四組決策機制讀取同一份鎖定之研究報告，差異僅在決策方式。}
\label{fig:flow}
\end{figure}

\subsection{對照組}

\paragraph{單次判斷（\grp{A}）。}
讀取研究報告後直接做出一次決策，代表不使用任何集成或辯論之基準。

\paragraph{獨立投票（\grp{B}）。}
同一提示詞以較高溫度獨立取樣七次，採多數決；平手時取信心度最高者。此組對應自我一致性方法 \citep{wang2023selfconsistency}，用以檢驗辯論之效果是否僅來自增加呼叫次數。

\paragraph{固定立場辯論（\grp{C}）。}
看多與看空代理人辯論三輪，每輪兩次呼叫，最後由中立裁決者決策，共七次呼叫。

\subsection{實驗組：立場交換辯論}

立場交換辯論（\grp{D}）第二輪兩位代理人互換立場，第三輪換回原立場。與 \grp{C} 不同，\grp{D} 交換輪僅能看到自己第一輪之論述，不看對手內容；代理人須引用並反駁自己第一輪的原主張（系統核對引用確實出自該代理人第一輪之文字，至少 70\% 字詞連續相同）。信心變化由系統以兩輪自報信心之差計算；模型另行回報之變化值僅保存於稽核紀錄（v3-1001.1 執行中調整，見事前登記補充說明）。隔離對手論述是為了讓代理人反駁自己而非回應對手，屬於處理的一部分；主要估計量為此自我反駁機制相對 \grp{C} 的效果。若要區分立場互換與資訊隔離各自的貢獻，需另設消融組。

裁決者須逐一核對雙方引用之證據是否確實支持其主張，未有證據支持之主張不予採計，並綜合雙方論點在買進與賣出間擇一（新版不得持有），所選方向須與其預期報酬之正負一致。長期記憶僅提供予辯論組之裁決者，且僅包含分析日前已到期之同股票案例。

\subsection{實驗設計}

\grp{B}、\grp{C}、\grp{D} 之決策呼叫次數對齊（各 7 次），\grp{A} 為單次判斷之基準（表~\ref{tab:arms}），以排除「呼叫次數較多因而表現較佳」之解釋；呼叫次數相同不代表 token 或運算量相同，辯論組因讀取歷史論述，輸入較長。

\begin{table}[H]
\centering\small
\caption{四組決策機制與模型呼叫次數}
\label{tab:arms}
\begin{tabularx}{\textwidth}{@{}c l c Y@{}}
\toprule
組別 & 名稱 & 呼叫次數 & 流程 \\
\midrule
\grp{A} & 單次判斷 & 1 & 讀取研究報告後做出一次決策 \\
\grp{B} & 獨立投票 & 7 & 七次獨立取樣，多數決 \\
\grp{C} & 固定立場辯論 & 7 & 看多與看空辯論三輪，再由裁決者決策 \\
\grp{D} & 立場交換辯論 & 7 & 第二輪互換立場但僅見自身第一輪論述，第三輪換回；\grp{C} 第二輪可見雙方第一輪論述 \\
\bottomrule
\end{tabularx}
\end{table}

每次決策輸出 60 個交易日之預期報酬、信心度、理由、引用證據與風險。舊版（v3-0929.1）之買進、持有或賣出由預期報酬相對於中性帶（歷史 60 日報酬標準差之一半）推導。新版季度 v3-1002.3／月度 v3-1002.4 改由模型直接選擇：買進即做多、賣出即放空、持有即不進場；\grp{C} 與 \grp{D} 之裁決者須綜合雙方論點、在買進與賣出間擇一，不得持有。中性帶僅保留為診斷用之推導行動。新版另將股票代號與公司名稱匿名化，以降低模型憑記憶辨識標的之前視風險；惟財報數值與日期仍可見，匿名化無法完全排除辨識。新版不要求模型於敘述中重述數字，而以證據代稱指稱事實（每位讀者皆收到附精確數值之證據清單）；敘述中若仍出現數字，程式比對證據數值後僅記錄無來源支持者，不修改文字。新舊規則產生的結果不可合併。

研究樣本於執行前以批次方式事前登記，僅登記後建立之實驗計入正式結果。分析日另依年份分為訓練（2021--2023）、驗證（2024）與測試（2025）三段。

正式 RQ1 使用 qwen3:32b（Ollama 0.35.0，部署於 RTX 5090 租用主機），溫度 0.2（\grp{B} 投票 0.8）、固定隨機種子、上下文 8192 tokens、單次輸出上限 2048 tokens，每次只送出一個呼叫。事前登記於 2026 年 10 月 2 日、任何案例執行前完成，協議雜湊為 \texttt{b1da80ee}，涵蓋所有設定、決定實驗結果之程式原始碼（資料、提示詞、驗證、回測等九個模組），並另行凍結模型權重之 digest；程式或權重若有改變，系統即拒絕排入或續跑，以確保 180 案皆在同一規則下產生。新版之前曾以相同 180 個資料集多次試跑與修正規則（作廢之登記與原因皆記錄於事前登記文件），且前導研究已看過同期結果，故本研究定位為「規則鎖定後對既有樣本之回溯比較」，而非全新之樣本外測試。

\subsection{績效衡量與統計檢定}

主要推論事前指定為：\grp{D} 對 \grp{C}、60 個交易日、扣除 Corwin--Schultz 成本、守門員之前的候選決策、全部案例。兩個共同主要端點為方向準確率（配對 McNemar 檢定 \citep{mcnemar1947}）與年化 Sharpe 比率差（Ledoit--Wolf 檢定 \citep{ledoit2008sharpe}）；Holm 校正 \citep{holm1979} 只在這兩者間進行，兩者差異皆為正且校正後 $p<0.05$，才判定 \grp{D} 優於 \grp{C}。其餘比較（\grp{D} 對 \grp{A}、\grp{B}，Jobson--Korkie 檢定 \citep{jobson1981sharpe}，30 與 90 日，零成本，守門員後決策）皆為敏感性分析。

持有視為棄權，另以覆蓋率（下注比例）、條件準確率（有下注時之準確率）與持有機會成本呈現；因 \grp{C}、\grp{D} 之裁決者必須下注而 \grp{A}、\grp{B} 可持有，與 \grp{A}、\grp{B} 之比較另以「兩組皆下注」之配對案例呈現。180 案共享 20 個分析日，故另報告以日期為群聚之信賴區間，並排除 37 個低新聞密度案例重算主要比較。所有結論等 180 個登記案例全數完成後才產生，不因中途結果提前停止。

RQ1 以 60 日為事前指定之主要期間，於同一模型內比較四組；30 與 90 個交易日為期間穩健性檢查，檢定與 Holm 校正於各期間內進行。RQ2 與 RQ3 沿用相同之比較，分別改變模型與分析頻率；RQ2 已以探索性分析完成（第~\ref{sec:exploratory} 節），RQ3 屬後續研究。

%% =====================================================================
\section{系統設計與實作}
\label{sec:system}
%% =====================================================================

本章說明實際完成並用於實驗之系統。系統核心不是單純線性串接代理人，而是由 Coordinator 控制任務、資料邊界與記憶權限，使四組決策機制只在報告生成之後才分流。

\subsection{系統架構}

使用者輸入股票代號、分析日期與任務設定後，Coordinator 先固定協議版本、模型與資料時間邊界，再分派研究任務給情緒面、基本面、技術面與總經代理人。各代理人只能讀取分析日以前已公開之證據，每筆證據有識別碼與可取得時間。研究結果合併為中立研究報告與證據登錄，報告不含買進、持有或賣出結論，鎖定後以內容雜湊固定；其後才依實驗組別進入單次判斷、獨立投票、固定立場辯論或立場交換辯論。守門員依引用有效率、資料覆蓋、波動度與基準率歷史對候選決策做最後檢查，回測模組再以真實價格計算結果。整體流程如圖~\ref{fig:flow}。

系統以網頁研究台操作（圖~\ref{fig:ui}）：使用者在同一頁面完成資料準備、實驗設定、Agent 執行與統計，資料流向圖即時顯示每個案例從資料來源、研究代理人、中立報告、四組決策到守門員與回測的進度，點選節點即可查看該步驟之輸入、模型輸出與稽核紀錄。

\begin{figure}[H]
\centering
\includegraphics[width=\textwidth]{figures/ui-pipeline.png}
\caption{網頁研究台之資料流向圖（一個已完成案例）。左至右依序為資料來源、資料代理人、資料快照、四個研究代理人、中立報告、四組決策機制、守門員、回測與統計。}
\label{fig:ui}
\end{figure}

\subsection{功能模組}

\begin{table}[H]
\centering\small
\caption{系統功能模組}
\label{tab:modules}
\begin{tabularx}{\textwidth}{@{}l Y Y@{}}
\toprule
模組 & 主要功能 & 輸出 \\
\midrule
Coordinator & 依股票與分析日派發任務，固定協議版本與模型，限制各代理人可讀取之資料 & 任務設定與流程狀態 \\
資料代理人 & 建立四個面向之分析日前快照，以內容雜湊為識別碼 & 不可變資料快照 \\
研究代理人 & 情緒面、基本面、技術面、總經面各自摘要；輸出須引用已提供之證據，驗證失敗時改用可稽核之來源摘錄 & 中立研究報告、證據登錄 \\
短期工作記憶 & 保存每個案例之全部呼叫紀錄（提示雜湊、種子、用量）、辯論與裁決紀錄 & 執行軌跡 \\
長期決策記憶 & 保存同一檔股票分析日前已到期案例之方向與報酬，僅提供辯論組（\grp{C}、\grp{D}）之裁決者 & 歷史案例摘要 \\
決策機制層 & 依組別執行單次判斷、七次獨立投票、固定立場辯論或立場交換辯論 & 候選決策 \\
裁決者 & 讀取報告、雙方辯論紀錄與記憶，逐項核對雙方引用之證據是否確實支持其主張，未有證據支持之主張不予採計 & 裁決方向、預期報酬、信心與理由 \\
守門員 & 檢查引用有效率、四面向覆蓋、年化波動度與基準率歷史 & 通過、降為持有或不交易 \\
回測模組 & 開盤對開盤計算 30、60、90 日報酬，扣除價差成本，並將結果寫回長期記憶 & 逐案例回測、彙總指標 \\
\bottomrule
\end{tabularx}
\end{table}

\subsection{立場交換辯論流程}

辯論固定為三輪，不設提前收斂偵測，使 \grp{C}、\grp{D} 兩辯論組之流程一致，並與 \grp{B} 之決策呼叫次數相同（\grp{A} 為單次判斷，表~\ref{tab:arms}）。表~\ref{tab:rounds} 為立場交換組（\grp{D}）之三輪流程。固定立場組（\grp{C}）輪次相同，但有兩處不同：第二輪不換邊，且第二輪可見雙方第一輪論述；\grp{D} 交換輪只見自己第一輪論述（見第~\ref{sec:methods} 節）。

\begin{table}[H]
\centering\small
\caption{三輪立場交換辯論流程}
\label{tab:rounds}
\begin{tabularx}{\textwidth}{@{}l Y Y@{}}
\toprule
輪次 & 內容 & 目的 \\
\midrule
第 1 輪 & 看多與看空代理人平行提出初始論點，引用證據，並指出對方論點可能忽略之缺口或假設 & 建立正反兩方初始立場，避免只呈現單方面敘事 \\
第 2 輪 & 雙方交換立場：原本看多者改為論證看空，原本看空者改為論證看多。只能看到自己第一輪之論述，須引用自己第一輪要推翻之主張（系統核對出處）並說明其為何不成立；信心變化由系統以兩輪信心之差計算 & 迫使代理人檢驗自身論點之弱點，避免僅改寫對方說法 \\
第 3 輪 & 雙方換回原始立場，整合前兩輪之論點與反駁，並列出仍無法確定之風險與資料缺口 & 產生整合後之最終論點與不確定性揭露 \\
裁決 & 裁決者讀取整份辯論紀錄、報告與記憶後輸出決策 & 逐項核對證據，避免以辯論聲量決定 \\
\bottomrule
\end{tabularx}
\end{table}

\subsection{驗證與可審查性}

為使每個決策可回溯，模型輸出寫入前皆經驗證，且每次呼叫皆留下稽核紀錄：

\begin{itemize}
  \item 提示詞中之證據以短代稱（E1、E2、$\ldots$）呈現，回答後由程式對映回真實識別碼，避免模型抄寫長識別碼出錯，同時不使識別碼中之申報編號洩漏公司身分。
  \item 引用之證據必須存在；引用不存在之證據或格式錯誤時，以明確訊息要求重試，最多三次。
  \item 數字檢查：新版要求代理人以證據代稱指稱事實、不重述數字；敘述中仍出現之數字若無任何證據數值支持（期間名稱、日期、來源值之四捨五入或截位不計），僅記錄並依組別報告比例，不重試、不修改文字。
  \item 凍結與身分核對：每一步推理前後皆核對模型權重之 digest，無法取得時暫停並可續跑，確認不同時則該案永久作廢；程式原始碼之雜湊納入協議雜湊，改動即拒絕續跑。連線故障連續發生時暫停派發新案例，並於停滯過久時自動關閉 GPU 主機。
  \item 每個案例可匯出完整紀錄：協議設定、中立報告、證據登錄、執行軌跡、記憶、決策、逐案例回測與彙總指標，以及辯論逐字稿。
\end{itemize}

\subsection{實作工具}

\begin{table}[H]
\centering\small
\caption{實作工具}
\label{tab:tools}
\begin{tabularx}{\textwidth}{@{}l l Y@{}}
\toprule
項目 & 工具 & 用途 \\
\midrule
流程控制 & LangGraph & 管理協調、研究、決策、守門與回測之狀態流 \\
模型呼叫 & Ollama、LiteLLM & 本地部署開源模型（部署於 GPU 租用主機）與雲端模型之統一呼叫 \\
資料 & yfinance、SEC EDGAR、FRED／ALFRED、FNSPID、Alpha Vantage & 行情、財報、總經與新聞 \\
情緒分析 & FinBERT & 逐則標題之正負面評分 \\
儲存與介面 & SQLite、網頁介面與命令列 & 保存案例、事前登記與匯出 \\
\bottomrule
\end{tabularx}
\end{table}

\subsection{實作挑戰與對策}

實作過程中遇到之主要問題與對策如下：分析日以後之資料不得進入輸入，故財報以申報日而非會計期末篩選、總經採當時已發布之版本；模型輸出格式不一致與引用錯誤，以結構化輸出、短代稱與重試機制處理；模型常自行計算或改寫數字，新版改以證據代稱取代重述數字，並記錄未匹配來源之數字；辯論文字過長，以輸出長度上限與欄位長度限制處理；證據不足時系統不應強行給出方向，由守門員將買進或賣出降為持有，或在基準率歷史不足時輸出不交易。守門員之規則為：年化波動度高於 0.80 時，買進或賣出降為持有；基準率歷史窗口不足時輸出不交易；引用有效率低於 80\% 或四面向未齊全時，依缺資料政策記錄，本研究採允許決策並標示；同股票歷史準確率低於 50\%（至少五筆）僅作警示，不改變決策。

%% =====================================================================
\section{實驗結果}
\label{sec:results}
%% =====================================================================

本節報告正式 RQ1（v3-1002.3）之結果、案例示範與探索性分析。

\subsection{正式 RQ1：自我反駁機制與固定立場辯論之比較}
\label{sec:formalresults}

180 個事前登記案例全數完成，無因模型權重不符而作廢、研究代理人降級或重複之案例；執行期間因網路短暫中斷而暫停之案例皆於恢復連線後續跑完成，未重新產生已完成之步驟。主要分析設定：候選決策層、60 個交易日、扣除 Corwin--Schultz 交易成本、四組等權重。

\begin{table}[H]
\centering\small
\caption{正式 RQ1 主要結果（qwen3:32b，v3-1002.3，60 日，扣除交易成本）}
\label{tab:formalmain}
\begin{tabular}{@{}lcccccc@{}}
\toprule
組別 & 覆蓋率 & 條件準確率 & 持有機會成本 & 總報酬 & 最大回撤 & 年化 Sharpe \\
\midrule
\grp{A}\ 單次判斷 & 78.3\% & 56.0\% & 7.51 & 1.85\% & $-1.81\%$ & 0.44 \\
\grp{B}\ 獨立投票 & 79.4\% & 55.9\% & 6.66 & 1.95\% & $-1.83\%$ & 0.46 \\
\grp{C}\ 固定立場辯論 & 100\% & 58.3\% & -- & 2.51\% & $-1.47\%$ & 0.59 \\
\grp{D}\ 立場交換辯論 & 100\% & 55.0\% & -- & 2.11\% & $-2.01\%$ & 0.54 \\
\midrule
買進持有（相同期間） & -- & -- & -- & 5.60\% & $-1.59\%$ & 1.01 \\
\bottomrule
\end{tabular}
\par\smallskip
\raggedright\footnotesize 組合將每個「股票×分析日」案例視為固定權重（1/180）之部位，僅於該案 60 日持有期間依決策做多或放空，其餘時間與持有之案例以零報酬現金計，各案持有期可能重疊；買進持有基準採相同配置、於每案相同期間做多該股票，並非九檔股票全期間持有。\grp{C}、\grp{D} 之裁決者必須下注，覆蓋率恆為 100\%；守門員可能因波動度將部分決策降為持有，該層結果列為敏感性分析。
\end{table}

\begin{table}[H]
\centering\footnotesize
\caption{正式 RQ1 主要推論與配對比較（60 日）}
\label{tab:formaltests}
\begin{tabular}{@{}lcccccc@{}}
\toprule
比較 & 配對案例 & 不一致 & 準確率差 & McNemar $p$ & 年化 Sharpe 差 & Ledoit--Wolf $p$ \\
\midrule
\grp{D} 對 \grp{C}（主要） & 180 & 16 & $-3.3$ & 0.21 & $-0.05$ & 0.80 \\
\quad Holm 校正（兩端點） & & & & 0.42 & & 0.80 \\
\grp{D} 對 \grp{B}（敏感性） & 143 & 9 & $+0.7$ & 1.00 & $+0.08$ & 0.79 \\
\grp{D} 對 \grp{A}（敏感性） & 141 & 8 & $+1.4$ & 0.73 & $+0.10$ & 0.72 \\
\bottomrule
\end{tabular}
\par\smallskip
\raggedright\footnotesize \grp{D} 對 \grp{C} 以 180 案配對；與 \grp{A}、\grp{B} 之比較以兩組皆下注之案例計算。\grp{D} 對 \grp{C} 之 16 個不一致案例中，\grp{D} 正確 5 個、\grp{C} 正確 11 個；準確率差單位為百分點。
\end{table}

\paragraph{敏感性分析。} 以分析日為群聚之區塊自助法，\grp{D}$-$\grp{C} 方向準確率差之 95\% 區間為 $[-7.2, +0.6]$ 個百分點（區塊長度 1、2、4 個日期結果相近），Sharpe 差之 Ledoit--Wolf 區間為 $[-0.47, +0.36]$，皆包含 0。排除 37 個低新聞密度案例後（143 案），準確率差 $-2.8$ 個百分點（$p=0.42$）、Sharpe 差 $-0.03$（$p=0.90$）。30 日與 90 日之準確率差皆為 $-1.1$ 個百分點（$p=0.80$），Sharpe 差分別為 $+0.04$ 與 $-0.03$（$p=0.93$、$0.88$）。不計成本及守門員後決策之結果與主要分析幾乎相同（守門員僅因 INTC 2024-09-30 年化波動度 83.6\% 將四組之賣出降為持有）。各敏感性分析皆未提供 \grp{D} 優於 \grp{C} 之顯著證據，但部分點估計方向不同（如 30 日 Sharpe 差為正）。

\paragraph{決策分布與文字品質。} 180 案中 129 案（72\%）四組決策完全相同；\grp{C} 與 \grp{D} 僅 16 案（9\%）方向不同，主要比較之檢定力集中於這些案例。以候選決策計，四組皆以買進為主（\grp{A} 買進 108、持有 39、賣出 33 案；\grp{B} 109、37、34 案；\grp{C} 買進 145、賣出 35 案；\grp{D} 買進 137、賣出 43 案），\grp{D} 之賣出多於 \grp{C}，\grp{A}、\grp{B} 出手率約 79\%。\grp{A} 與 \grp{B} 之分布幾乎相同，顯示七次投票之多數決多與單次判斷一致。\grp{D} 交換輪引用之證據 ID 新穎率平均 0.34，\grp{C} 為 0.05（參照輪不同，僅為描述）。含未匹配來源數字之呼叫比例：決策呼叫整體 6.9\%（\grp{A} 4.4\%、\grp{B} 6.2\%、\grp{C} 5.3\%、\grp{D} 9.7\%），研究代理人技術面 47.8\%、情緒面 15.6\%、總經面 0\%；技術面多為照抄指標值之變體寫法，此指標僅代表與來源數值不相符，不等於事實錯誤率。

\subsection{案例示範：自我反駁如何改變裁決}
\label{sec:casestudy}

為說明兩組流程之實際差異，以下呈現 \grp{C} 與 \grp{D} 方向不同之兩個案例：一個 \grp{D} 判斷正確，一個判斷錯誤。兩案之研究報告、證據與模型完全相同，差別僅在辯論流程。代理人之原文為英文，以下為摘譯，括號內為其引用之證據代稱。

\paragraph{案例一：INTC 2024-03-31（\grp{D} 正確）。} 研究報告顯示新聞情緒偏正面（FinBERT 平均 0.25，正面比例 67\%）、近 20 日報酬為正，但營收、淨利與營業現金流皆較去年同期大幅衰退。表~\ref{tab:case} 為兩組之辯論與裁決。

\begin{table}[H]
\centering\small
\caption{案例一：INTC 2024-03-31 之辯論與裁決（60 日實際報酬 $-30.5\%$）}
\label{tab:case}
\begin{tabularx}{\textwidth}{@{}l l Y@{}}
\toprule
組別／輪次 & 立場與預測 & 論點摘要 \\
\midrule
\grp{C} 第 1--3 輪 & 看多 $+8\%\rightarrow+7\%$；看空 $-12\%\rightarrow-9\%$ & 看多方三輪幾乎重複同一論點：「情緒正面（E4、E5）與短期動能（E11）支持看多，雖然基本面衰退」；看空方亦三輪重複財報衰退。雙方未回應對方之核心證據。 \\
\grp{C} 裁決 & 買進，$+5\%$ & 採納看多方之「短期動能與情緒」論點。 \\
\midrule
\grp{D} 第 1 輪 & 看多 $+8\%$；看空 $-10\%$ & 與 \grp{C} 第 1 輪相同之兩種論點。 \\
\grp{D} 第 2 輪（交換） & 原看多者改看空 $-7\%$ & 引用並反駁自己第 1 輪之主張「正面情緒與短期報酬顯示上漲動能」，改稱「這些訊號被營收、淨利與現金流之大幅衰退（E1--E3）所抵銷」。 \\
\grp{D} 第 3 輪 & 看多 $+6.5\%$；看空 $-8\%$ & 看空方整合交換輪之論點：基本面惡化之重要性高於短期情緒。 \\
\grp{D} 裁決 & 賣出，$-7\%$ & 「營收、淨利與營業現金流之年減顯示基本面惡化；情緒與短期報酬雖正面，但不足以抵銷。」 \\
\bottomrule
\end{tabularx}
\end{table}

此案中，交換輪迫使原看多之代理人正視自己第一輪輕描淡寫之基本面衰退，裁決者隨之轉為賣出；60 日後股價下跌 30.5\%。\grp{A}、\grp{B} 亦判斷賣出，\grp{C} 是四組中唯一判斷錯誤者。

\paragraph{案例二：AAPL 2023-03-31（\grp{D} 錯誤）。} 同樣的機制也可能帶來反效果。此案技術面強勁（近 20 日上漲 11.3\%），但營收與淨利年減。\grp{C} 之裁決者採納技術動能判斷買進；\grp{D} 交換輪中原看多者改為看空，強調基本面衰退與 83.6\% 之負債比，裁決者隨之判斷賣出。60 日後股價上漲 15.3\%，\grp{D} 錯誤。

\paragraph{分歧案例之整體模式。} 180 案中 \grp{C} 與 \grp{D} 方向不同之 16 案裡，\grp{D} 有 12 案判斷賣出，正確 5 案、錯誤 11 案。兩個案例中皆可觀察到交換輪強調反面證據（多為基本面衰退），且最終決策轉為看空，此模式與自我反駁之設計目的相符；但少數案例不足以確認其因果機制。可能的解釋是：本研究期間以上漲為主，基本面疲弱而股價續漲之情形較多，較常看空反而不利。立場交換改變了辯論內容與部分決策，但未帶來更準確的結果。

\paragraph{對 RQ1 之回答。} 在事前登記之 180 個季度案例中，自我反駁式立場交換辯論（\grp{D}）未優於固定立場辯論（\grp{C}）：兩個共同主要端點之點估計皆略偏向 \grp{C}（準確率 $-3.3$ 個百分點、Sharpe $-0.05$），但皆不顯著，各敏感性分析亦未提供 \grp{D} 較佳之顯著證據。要求代理人站到對面反駁自己後，可觀察到其引用不同證據、決策較常轉為看空，但未轉化為更準確或風險調整後更佳之決策。四組皆未勝過同期買進持有；辯論兩組之 Sharpe 雖高於 \grp{A}、\grp{B}，但兩組必須下注，在以上漲為主之樣本期間較少持有，差異不能歸因於辯論本身。

\subsection{探索性分析：跨模型複製與基準比較}
\label{sec:exploratory}

以下分析皆於主要結果揭曉後進行，未列入事前登記，依 \citet{wagenmakers2012agenda} 與 \citet{nosek2018prereg} 之區分屬探索性分析，僅用於產生後續研究之假說，不作為確認性證據。

\paragraph{跨模型複製。} 以 Gemini 3 Flash 依相同流程、相同 180 案與相同匿名化設定重跑四組（表~\ref{tab:flash}）。結果與 qwen3:32b 一致：\grp{D} 之準確率仍略低於 \grp{C}，\grp{C}、\grp{D} 方向不同之案例僅 9 案（\grp{D} 單獨正確 3 案、\grp{C} 單獨正確 6 案，McNemar $p=0.51$）。兩模型之決策高度相似，\grp{C}、\grp{D} 分別有 165 與 162 案方向相同。

\begin{table}[H]
\centering\small
\caption{跨模型比較：方向準確率（60 日、候選決策、有出手之案例）}
\label{tab:flash}
\begin{tabular}{@{}lcc@{}}
\toprule
組別 & qwen3:32b & Gemini 3 Flash \\
\midrule
\grp{A} 單次判斷 & 56.0\% & 57.6\% \\
\grp{B} 獨立投票 & 55.9\% & 56.7\% \\
\grp{C} 固定立場辯論 & 58.3\% & 57.8\% \\
\grp{D} 立場交換辯論 & 55.0\% & 56.1\% \\
\midrule
一律看多（事後參照） & \multicolumn{2}{c}{62.2\%} \\
\bottomrule
\end{tabular}
\end{table}

\paragraph{與一律看多之比較。} 180 案中 112 案於 60 日內上漲（62.2\%）。兩模型四組之準確率皆低於一律看多；68 個下跌案例中，各組僅判斷正確 14--15 案（約 21\%）。此一參照需在事後才能得知，但顯示現有輸入對下跌幾乎沒有辨識力。

\paragraph{被否證之解釋。} 本研究原推測 \grp{D} 較常看空，因而在多頭期吃虧、在空頭期應較佳。將 20 個分析日依九檔平均 60 日報酬分為 6 個空頭季與 14 個多頭季後，\grp{D} 在空頭季之準確率亦未高於 \grp{C}（qwen 24.1\% 對 25.9\%；Flash 25.9\% 對 27.8\%），此推測不成立。

\paragraph{後續假說。} 綜合上述，推理結構（立場交換）與模型能力之改變皆未改善決策，而四組均低於一律看多，可能之解釋為限制來自輸入資訊之預測力。此與部分文獻發現辯論可提升推理任務表現 \citep{du2023debate,liang2023divergent} 不同，可能原因在於本研究為高雜訊之預測任務，惟本研究未檢驗此差異。後續研究將保持流程不變、僅擴充資訊集（如跨標的聯動訊號），並依 \citet{nosek2018prereg} 之建議，於新期間或新標的之資料上事前登記檢驗，以避免在產生假說之同一批資料上驗證。

\subsection{研究限制}

第一，研究範圍為九檔美國大型股，結果未必能推廣至其他市場或小型股。第二，D 的自我反駁機制由立場互換、隔離對手與強制引用自身主張組成，本研究未以消融組區分各要素的個別貢獻。第三，標的匿名化僅隱藏代號與公司名稱，財報數值與日期仍可見，模型仍可能推知標的；且 2021--2025 年可能在模型訓練資料內，本研究以匿名化與記憶探測降低、但無法排除前視偏差。第四，本批為規則鎖定後對既有樣本之回溯比較，前導研究已看過同期結果，不宣稱全新之樣本外表現。第五，\grp{C}、\grp{D} 裁決者必須下注且可讀取已到期之歷史記憶，\grp{A}、\grp{B} 可持有且無記憶，故 \grp{D} 對 \grp{A}、\grp{B} 之差異混合多項設計差異，只有 \grp{D} 對 \grp{C} 為主要比較。第六，180 案共享 20 個日期，90 日預測期間又長於季度間隔，方向檢定另以日期群聚區間檢查。第七，新聞來源混用 FNSPID 與 Alpha Vantage，36 案僅有即時標題，新聞密度與年份、個股相關。本研究未對每個案例重複執行，無法評估同案重跑變異。

%% =====================================================================
\section{結論與未來研究}
\label{sec:conclusion}
%% =====================================================================

正式 RQ1（qwen3:32b，v3-1002.3，180 個事前登記案例）結果為：\grp{D} 自我反駁機制未優於 \grp{C} 固定立場辯論，方向準確率 55.0\% 對 58.3\%（$p=0.21$）、年化 Sharpe 0.54 對 0.59（$p=0.80$），經 Holm 校正後皆不顯著，30／90 日、零成本、守門員後與排除低新聞案例之敏感性分析亦未提供 D 較佳之顯著證據；四組均未勝過同期買進持有（Sharpe 1.01）。可觀察到自我反駁使代理人引用較多新證據、決策較常看空，但未帶來更好之決策。


RQ2 之跨模型比較已以 Gemini 3 Flash 於同一批 180 案完成，惟未事前登記，僅作探索性分析：其結論與 qwen3:32b 相同，且兩模型各組準確率皆低於一律看多，顯示限制可能來自輸入資訊。RQ3 之月度擴充（v3-1002.4，540 案）尚未執行。未來研究將於新期間或新標的上事前登記，檢驗擴充資訊集（如跨標的聯動訊號）之效果。若要區分立場互換與資訊隔離各自的貢獻，另須新增消融組（例如僅讀自身歷史但不換立場）並重新事前登記。

\vspace{0.8cm}
\noindent\small 本文為國立陽明交通大學管理科學系專題研究之成果，僅供學術研究與教學使用，不構成任何投資建議。
\normalsize

\clearpage
\bibliographystyle{chicago}
\bibliography{references}

\clearpage
\end{CJK}
\end{document}
